\documentclass[10pt,a4paper]{article}
\usepackage[margin=2.2cm]{geometry}
\usepackage{fontspec}
\setmonofont{DejaVu Sans Mono}[Scale=0.84]
\newfontfamily\unifont{DejaVu Serif}[Scale=0.9]
\usepackage{amsmath,amssymb,mathtools}
\setcounter{MaxMatrixCols}{20}
\usepackage{longtable,booktabs,array}
\usepackage{fancyvrb}
\usepackage[dvipsnames]{xcolor}
\usepackage[colorlinks=true,linkcolor=NavyBlue,urlcolor=NavyBlue]{hyperref}
\usepackage{microtype}
\usepackage{needspace}
\usepackage{tocloft}\setlength{\cftbeforesecskip}{0.45em}
\setlength{\parindent}{0pt}\setlength{\parskip}{0.55em}
\setlength{\emergencystretch}{8em}
\makeatletter
\renewcommand\section{\@startsection{section}{1}{\z@}{-3.5ex \@plus -1ex \@minus -.2ex}{2.3ex \@plus.2ex}{\normalfont\Large\bfseries\raggedright}}
\renewcommand\subsection{\@startsection{subsection}{2}{\z@}{-3.25ex\@plus -1ex \@minus -.2ex}{1.5ex \@plus .2ex}{\normalfont\large\bfseries\raggedright}}
\renewcommand\subsubsection{\@startsection{subsubsection}{3}{\z@}{-3.25ex\@plus -1ex \@minus -.2ex}{1.5ex \@plus .2ex}{\normalfont\normalsize\bfseries\raggedright}}
\makeatother
\renewcommand{\arraystretch}{1.15}
\allowdisplaybreaks
\begin{document}
\title{The world brane of 7+1 spacetime: symbols and their definitions\\[0.4em]\large A dictionary of every symbol used in the document on the world brane, with its defining equation and its name in the notebook 7+1.nb}
\author{Repository: github.com/once-ere/7plus1 — prepared for Nash\\ Notebook: 7+1.nb (Wolfram Language 15.0.1)\\ Document generated by tools/make\_docs.py on 2026-10-28}
\date{}
\maketitle
\begin{abstract}
For each symbol this dictionary gives the mathematical notation, the name used in the Mathematica notebook 7+1.nb (where there is one; the names of section 25 begin with \texttt{tb}, those of section 26 with \texttt{rd}; "—" means that the notebook has no separate name), the defining equation and its meaning, grouped by topic in the order of the document. Longer defining equations follow each table as numbered displayed formulas. Section 14 lists the symbols that look alike, section 15 the Mathematica names. Every equation needed to define a symbol is written in this document.
\end{abstract}
\tableofcontents
\Needspace{12\baselineskip}\section*{1. Conventions}\addcontentsline{toc}{section}{1. Conventions}
\begin{itemize}
\item Units: $\hbar=c=1$; energies, masses and inverse lengths in GeV; $\hbar c=1.973269804\times10^{-16}$ GeV m (Mathematica \texttt{tbHbarCm}).
\item Coordinates $X=(x_0,x_1,\dots,x_7)$ (Mathematica \texttt{X} $=$ \texttt{\{x0, …, x7\}}): $x_0$ is time (the direction of the octonion unit $e_0=1$), $x_1,x_2,x_3$ ordinary space, $x_4,\dots,x_7$ the four auxiliary space-like coordinates $\vec y$; on the brane $x=(x_0,\dots,x_3)$.
\item Indices: $A,B=0,\dots,7$ (all directions); $\mu,\nu=0,\dots,3$ (along the brane); $k,l=4,\dots,7$ (transverse to it); $i,j=1,2,3$ (ordinary space); $a=1,\dots,4$ (the transverse spinor index). A prime is $d/dx_0$ (or $d/dy$ for warped metrics), a dot $d/ds$.
\item Curvature (mostly-minus signature): $R^\rho{}_{\sigma\mu\nu}=\partial_\mu\Gamma^\rho{}_{\nu\sigma}-\partial_\nu\Gamma^\rho{}_{\mu\sigma}+\Gamma^\rho{}_{\mu\lambda}\Gamma^\lambda{}_{\nu\sigma}-\Gamma^\rho{}_{\nu\lambda}\Gamma^\lambda{}_{\mu\sigma}$, $R_{\sigma\nu}=R^\rho{}_{\sigma\rho\nu}$, $R=g^{\sigma\nu}R_{\sigma\nu}$, $G^A{}_B=R^A{}_B-\frac12\delta^A_BR$; $G^A{}_B-\Lambda\delta^A_B=\kappa T^A{}_B$.
\item Status labels used in the document: an ESTIMATE is an order of magnitude with numerical factors omitted; everything else is an exact consequence of the definitions, checked by sections 25–26 of 7+1.nb and by tools/ref\_torus\_brane.py and tools/ref\_radion.py.
\end{itemize}
\textbf{The signature: $\{7,1,0\}$ and $\mathrm{diag}(+1,-1,\dots,-1)$.} The 7+1 spacetime has one time-like dimension ($x_0$, the direction of the octonion unit $e_0=1$) and seven space-like dimensions ($x_1,x_2,x_3$ ordinary space, $x_4,\dots,x_7$ auxiliary space). Two different objects describe this, and neither may be replaced by the other. (i) The tensor package xAct, used in section 12 of 7+1.nb exactly as in the author’s notebook Einstein-Lovelock-7+1-Nash.nb, declares the metric with \texttt{DefMetric[\{7, 1, 0\}, g71[-a, -b], CD, …]}: the list $\{7,1,0\}$ declares seven space-like dimensions, one time-like dimension and no degenerate dimension — the numbers of entries $+1$ and $-1$ of $\eta_{71}=\mathrm{diag}(-1,+1,\dots,+1)$ (\texttt{η71}), the same spacetime written in the sign convention of 4+4.nb (space-like $=+1$, time-like $=-1$). (ii) The matrix $\eta=-\eta_{71}=\mathrm{diag}(+1,-1,\dots,-1)$ (\texttt{η17}) is simply the squares of the octonion basis: $e_0^2=+1$ for the unit $e_0=1$ (time) and $e_i^2=-1$ for the seven imaginary units $e_1,\dots,e_7$ (space); more generally it is the real part of the product of two basis units, and it gives the quadratic form of the octonions:
\begin{gather*}
\eta_{AB}=\mathrm{Re}(e_Ae_B),\qquad e_0^2=+1,\qquad e_i^2=-1\ \ (i=1,\dots,7),\\ \mathrm{Re}(x\,x)=x_0^2-x_1^2-\cdots-x_7^2=\eta_{AB}\,x^Ax^B\qquad (x=x^Ae_A\in\mathbb O).
\end{gather*}
$\eta$ gives the signs of the Clifford relation $\Gamma_A\Gamma_B+\Gamma_B\Gamma_A=2\,\eta_{AB}\,1_{16}$ of the real 16×16 Dirac matrices: time squares to $+1_{16}$ and each of the seven space directions to $-1_{16}$, exactly like the octonion units (in the notation $Cl(p,q)$, with $p$ generators squaring to $+1$ and $q$ to $-1$, the algebra $Cl(1,7)\cong\mathbb R(16)$). It also gives the signs of the metric components: $g_{00}>0$ along the time-like direction, negative squares along the space-like directions. Every component computed by 7+1.nb follows from the component matrix of the metric; the declared list fixes for xAct the sign of the determinant of the metric, \texttt{SignDetOfMetric[g71]} $=-1$, which is the same in both sign conventions, since $\det \eta=\det \eta_{71}=-1$. 7+1.nb checks \texttt{SignatureOfMetric[g71]} $=\{7,1,0\}$, \texttt{SignDetOfMetric[g71]} $=-1$, $\mathrm{Re}(e_Ae_B)=\eta_{AB}$ and $\mathrm{Re}(x\,x)=\eta_{AB}x^Ax^B$.

\Needspace{12\baselineskip}\section*{2. The bulk, the torus and the brane}\addcontentsline{toc}{section}{2. The bulk, the torus and the brane}
{\footnotesize
\begin{longtable}{>{\raggedright\arraybackslash}p{\dimexpr 0.1700\linewidth-2\tabcolsep\relax}>{\raggedright\arraybackslash}p{\dimexpr 0.2200\linewidth-2\tabcolsep\relax}>{\raggedright\arraybackslash}p{\dimexpr 0.3100\linewidth-2\tabcolsep\relax}>{\raggedright\arraybackslash}p{\dimexpr 0.3000\linewidth-2\tabcolsep\relax}}
\toprule
\textbf{symbol} & \textbf{Mathematica} & \textbf{definition} & \textbf{meaning} \\
\midrule\endfirsthead
\toprule
\textbf{symbol} & \textbf{Mathematica} & \textbf{definition} & \textbf{meaning} \\
\midrule\endhead
\endfoot
\bottomrule\endlastfoot
$\vec x$, $\vec y$ & \texttt{\{x1, x2, x3\}}, \texttt{tbY} & $\vec y=(x_4,\dots,x_7)$ & ordinary space; the auxiliary coordinates \\
$L_c$, $n_k$ & \texttt{tbLc} & (2.1) & coordinate period; integers \\
$T^4$ & — & (2.1) & the auxiliary torus \\
$M$ & — & (2.1) & the bulk \\
$b$, $b_0$ & \texttt{b4[x0]}, \texttt{tbB} & (2.1) & size of the torus; its value today \\
$V_\perp$, $\ell_\perp$ & \texttt{tbV}, \texttt{rIIVp}, \texttt{tbLp} & (2.1) & proper volume and size of the torus \\
$W$ & \texttt{tbOnW} (the rule $\vec y\to0$) & (2.2) & the world brane \\
$X^A(x)$, $J^A{}_\mu$ & \texttt{tbJ} & (2.2) & embedding of the brane; its Jacobian \\
$x^\mu$ & — & $x^\mu=X^\mu$ & brane coordinates \\
$g_{AB}$ & \texttt{gCos}, \texttt{gPoly} & section 3 of this dictionary & bulk metric \\
$h_{\mu\nu}$ & \texttt{tbInduced[g]} & (2.3) & induced metric of the brane \\
$e_0,\dots,e_3$ & — & (2.4) & basis of the quaternion subalgebra $\mathbb H$ of the octonions \\
$\epsilon_{ijk}$ & — & totally antisymmetric, $\epsilon_{123}=1$ & Levi-Civita symbol \\*
\end{longtable}}
\begin{align*}
x_k&\simeq x_k+L_cn_k,\qquad T^4=\mathbb R^4/(L_c\mathbb Z)^4,\qquad M=\mathbb R_{x_0}\times\mathbb R^3_{\vec x}\times T^4_{\vec y},&&(2.1) \\
b^8&=\det(-g_{kl}),\qquad V_\perp=b^4L_c^4,\qquad \ell_\perp=bL_c, \\
W&=\{X\in M:\ x_4=x_5=x_6=x_7=0\},&&(2.2) \\
X^A(x)&=(x_0,x_1,x_2,x_3,0,0,0,0),\qquad J^A{}_\mu=\frac{\partial X^A}{\partial x^\mu}, \\
h_{\mu\nu}&=g_{AB}J^A{}_\mu J^B{}_\nu=g_{\mu\nu}(x,\vec y=0),\qquad h=J^{\mathsf T}gJ,&&(2.3) \\
e_1e_2&=-e_3,\quad e_2e_3=-e_1,\quad e_3e_1=-e_2,\quad e_ie_j=-\delta_{ij}-\epsilon_{ijk}e_k,&&(2.4) \\
\mathrm{Re}(e_\mu e_\nu)&=\mathrm{diag}(1,-1,-1,-1).
\end{align*}
\Needspace{12\baselineskip}\section*{3. Metrics and the action}\addcontentsline{toc}{section}{3. Metrics and the action}
{\footnotesize
\begin{longtable}{>{\raggedright\arraybackslash}p{\dimexpr 0.1700\linewidth-2\tabcolsep\relax}>{\raggedright\arraybackslash}p{\dimexpr 0.2200\linewidth-2\tabcolsep\relax}>{\raggedright\arraybackslash}p{\dimexpr 0.3100\linewidth-2\tabcolsep\relax}>{\raggedright\arraybackslash}p{\dimexpr 0.3000\linewidth-2\tabcolsep\relax}}
\toprule
\textbf{symbol} & \textbf{Mathematica} & \textbf{definition} & \textbf{meaning} \\
\midrule\endfirsthead
\toprule
\textbf{symbol} & \textbf{Mathematica} & \textbf{definition} & \textbf{meaning} \\
\midrule\endhead
\endfoot
\bottomrule\endlastfoot
$n$, $a$, $b$ & \texttt{nP}, \texttt{aP}, \texttt{bP}; \texttt{a3[x0]}, \texttt{b4[x0]} & (3.1) & lapse, scale factors of ordinary space and of the torus \\
$\tau$ & — & time since the singularity & time of the primordial regime \\
$p_1$, $q_1$ & \texttt{rIp1}, \texttt{rIq1} & (3.2) & Einstein Kasner exponents \\
$S$, $S_{\rm grav}$, $S_\psi$, $S_U$, $S_W$ & — & (3.3) & the action and its parts \\
$\kappa$, $G_8$, $M_8$ & \texttt{tbKap}, \texttt{rIIM8} & $\kappa=8\pi G_8=M_8^{-6}$ & eight-dimensional gravity \\
$\Lambda$ & \texttt{lamb} & (3.3) & bulk cosmological constant \\
$R$, $L_2$, $L_3$, $\alpha_2$, $\alpha_3$ & \texttt{alf[2]}, \texttt{alf[3]} & (3.4) & scalar curvature, Lovelock densities and couplings \\
$\psi$, $\bar\psi$ & \texttt{dirac16complex} & $\bar\psi=\psi^\dagger\Gamma_0$ & the 16-component bulk field and its conjugate \\
$\Gamma^A(X)$, $\Omega_A$ & — & frame Dirac matrices, spin connection & Dirac operator of the bulk \\
$m$ & — & mass & mass of dirac16complex \\
$S_{\rm self}$ & — & nine contact terms & self-couplings of dirac16complex \\
$U(b)$, $U_0$, $K_U$ & \texttt{uPot}, \texttt{u0B}, \texttt{kU} & (5.4) & stabilizing potential \\
$f_W$ & \texttt{tbfW}, \texttt{fWb} & (4.4) & brane tension \\
$\mathcal L_{\rm SM}$, $\mathcal L_{\rm portal}$ & — & sections 8, 9 of this dictionary & Standard Model; portals \\*
\end{longtable}}
\begin{align*}
ds^2&=n^2dx_0^2-a^2d\vec x^2-b^2d\vec y^2\qquad\text{(inhomogeneous example:)}&&(3.1) \\
n&=1+\tfrac{x_0x_4}5+\tfrac{x_0^2}{11},\qquad a=1+\tfrac{x_0}3+\tfrac{x_4^2}7,\qquad b=2-\tfrac{x_0x_4}4+\tfrac{x_0^2}9+\tfrac{x_4}{13}, \\
p_1&=\frac{1+2\sqrt2}7,\qquad q_1=\frac{2-3\sqrt2}{14},\qquad a\propto\tau^{p_1},\ b\propto\tau^{q_1},&&(3.2) \\
S&=S_{\rm grav}+S_\psi+S_U+S_W,\qquad S_{\rm grav}=-\frac1{2\kappa}\int d^8X\sqrt{|g|}\,(2\Lambda+R+\alpha_2L_2+\alpha_3L_3),&&(3.3) \\
S_W&=\int_Wd^4x\sqrt{|h|}\,(-f_W+\mathcal L_{\rm SM}+\mathcal L_{\rm portal}), \\
L_2&=R^2-4R_{AB}R^{AB}+R_{ABCD}R^{ABCD},&&(3.4) \\
L_3&=R^3-12RR_{AB}R^{AB}+16R^A{}_BR^B{}_CR^C{}_A+24R^{AC}R^{BD}R_{ABCD}+3RR_{ABCD}R^{ABCD} \\
&\quad-24R^{AB}R_A{}^{CDE}R_{BCDE}+4R^{AB}{}_{CD}R^{CD}{}_{EF}R^{EF}{}_{AB}-8R_A{}^C{}_B{}^DR_C{}^E{}_D{}^FR_E{}^A{}_F{}^B .
\end{align*}
\Needspace{12\baselineskip}\section*{4. How the brane sits in the bulk; matter on the brane}\addcontentsline{toc}{section}{4. How the brane sits in the bulk; matter on the brane}
{\footnotesize
\begin{longtable}{>{\raggedright\arraybackslash}p{\dimexpr 0.1700\linewidth-2\tabcolsep\relax}>{\raggedright\arraybackslash}p{\dimexpr 0.2200\linewidth-2\tabcolsep\relax}>{\raggedright\arraybackslash}p{\dimexpr 0.3100\linewidth-2\tabcolsep\relax}>{\raggedright\arraybackslash}p{\dimexpr 0.3000\linewidth-2\tabcolsep\relax}}
\toprule
\textbf{symbol} & \textbf{Mathematica} & \textbf{definition} & \textbf{meaning} \\
\midrule\endfirsthead
\toprule
\textbf{symbol} & \textbf{Mathematica} & \textbf{definition} & \textbf{meaning} \\
\midrule\endhead
\endfoot
\bottomrule\endlastfoot
$n_{(k)}$ & — & $b^{-1}\partial_k$ & unit normals of the brane \\
$K^{(k)}_{\mu\nu}$ & — & (4.1) & extrinsic curvature \\
$\{\lambda\,\mu\nu\}$ & \texttt{tbChrW} & (4.1) & Christoffel symbols \\
$s$ & — & proper time & parameter of a geodesic \\
$G_\perp$ & — & (4.2) & Green’s function of the four transverse directions \\
$\delta^4(\vec y)$, $F$ & \texttt{DiracDelta} & (4.3) & delta function of the torus; a test function \\
$\vec q_{\vec n}$ & \texttt{tbQ[n]} & $2\pi\vec n/L_c$ & wave vector of a mode \\
$T^W_{\mu\nu}$, $T^{AB}$ & — & (4.4) & brane energy–momentum; its eight-dimensional source \\
$\varepsilon$, $p$, $w$ & — & (4.4) & energy density, pressure, $w=p/\varepsilon$ \\
$p_a$, $p_b$ & \texttt{feStress[W]} & (4.5) & pressures along ordinary space and along the torus \\
$W(a,b)$ & argument of \texttt{feStress} & (4.5) & comoving energy of a unit coordinate cell \\
$H_a$, $H_b$ & \texttt{Ha}, \texttt{Hb} & $a'/a$, $b'/b$ & expansion rates \\
$\rho_4$, $p_4$ & \texttt{tbRho} & (4.5) & energy density and pressure of a brane fluid \\
$k_a$, $k_b$, $\omega$, $\theta$ & \texttt{tbAa}, \texttt{tbBb}, \texttt{tbth} & (4.6) & comoving momenta, energy and angle of a massless particle \\
$u_0,\dots,u_3$ & \texttt{tbu0} … \texttt{tbu3} & $h=\mathrm{diag}(1/u_0,-1/u_1,-1/u_2,-1/u_3)$ & components used in the tension check \\*
\end{longtable}}
\begin{align*}
K^{(k)}_{\mu\nu}&=\frac12(\mathcal L_{n_{(k)}}g)_{\mu\nu}=\frac1{2b}\partial_kg_{\mu\nu}\Big|_W,&&(4.1) \\
\{\lambda\,\mu\nu\}&=\frac12g^{\lambda\sigma}(\partial_\mu g_{\sigma\nu}+\partial_\nu g_{\sigma\mu}-\partial_\sigma g_{\mu\nu}), \\
\nabla_{\vec y}^2\frac1{|\vec y|^2}&=-4\pi^2\delta^4(\vec y),\qquad G_\perp=-\frac1{4\pi^2|\vec y|^2},&&(4.2) \\
\int_Wd^4x\sqrt{|h|}\,\mathcal L&=\int d^8X\sqrt{|g|}\,\frac{\delta^4(\vec y)}{b^4}\,\mathcal L,\qquad \delta^4(\vec y)=\frac1{L_c^4}\sum_{\vec n}e^{i\vec q_{\vec n}\cdot\vec y},&&(4.3) \\
T^W_{\mu\nu}&=\frac2{\sqrt{|h|}}\frac{\delta S_W}{\delta h^{\mu\nu}},\qquad T^{AB}=\frac{\delta^4(\vec y)}{b^4}T_W^{\mu\nu}\delta^A_\mu\delta^B_\nu,&&(4.4) \\
\mathcal L_W&=-f_W\ \Rightarrow\ T^W_{\mu\nu}=f_Wh_{\mu\nu},\quad \varepsilon=f_W,\quad p=-f_W,\quad w=-1, \\
\varepsilon&=\frac W{a^3b^4},\qquad p_a=-\frac{\partial_aW}{3a^2b^4},\qquad p_b=-\frac{\partial_bW}{4a^3b^3},&&(4.5) \\
&\text{tension: }W=\frac{a^3f_W}{L_c^4};\qquad\text{brane fluid: }W=\frac{a^3\rho_4}{L_c^4}, \\
W&=\omega=\sqrt{k_a^2/a^2+k_b^2/b^2},\qquad k_a/a=p\cos\theta,\quad k_b/b=p\sin\theta,&&(4.6) \\
&\text{isotropic in seven dimensions: }p_a=p_b=\frac\varepsilon7.
\end{align*}
\Needspace{12\baselineskip}\section*{5. Gravity seen from the brane; the radion}\addcontentsline{toc}{section}{5. Gravity seen from the brane; the radion}
{\footnotesize
\begin{longtable}{>{\raggedright\arraybackslash}p{\dimexpr 0.1700\linewidth-2\tabcolsep\relax}>{\raggedright\arraybackslash}p{\dimexpr 0.2200\linewidth-2\tabcolsep\relax}>{\raggedright\arraybackslash}p{\dimexpr 0.3100\linewidth-2\tabcolsep\relax}>{\raggedright\arraybackslash}p{\dimexpr 0.3000\linewidth-2\tabcolsep\relax}}
\toprule
\textbf{symbol} & \textbf{Mathematica} & \textbf{definition} & \textbf{meaning} \\
\midrule\endfirsthead
\toprule
\textbf{symbol} & \textbf{Mathematica} & \textbf{definition} & \textbf{meaning} \\
\midrule\endhead
\endfoot
\bottomrule\endlastfoot
$G_4$, $M_{\rm Pl}$ & \texttt{tbGN}, \texttt{tbMP}, \texttt{rIIMPl} & (5.1) & four-dimensional Newton constant; reduced Planck mass \\
$\Lambda_4$ & — & (5.1) & four-dimensional cosmological constant \\
$H_0$, $\Omega_m$, $\Omega_r$, $\Omega_\Lambda$, $t_0$ & \texttt{rIIH0g}, \texttt{rIIOL} & (5.2) & Hubble rate, density fractions, age today \\
$\rho_\Lambda$ & \texttt{dkRhoL} & $3H_0^2M_{\rm Pl}^2\Omega_\Lambda$ & observed dark-energy density \\
$\Phi_8$, $\Phi_{\rm N}$ & — & (5.3) & Newtonian potentials (eight-dimensional; on the torus) \\
$M$, $M'$, $r$ & \texttt{tbr} & — & two masses on the brane; their distance \\
$\beta$, $\phi$ & \texttt{rdBt}, \texttt{rdBe} & (5.4) & the radion; the canonical radion \\
$g_E$ & — & (5.4) & four-dimensional Einstein-frame metric \\
$m_r$ & — & 10 GeV & radion mass \\
$T^\mu{}_\mu$, $\mathcal L_{\phi W}$, $V_\phi$ & — & (5.4) & trace of the brane energy–momentum; radion coupling; radion potential \\
$\Gamma_{\rm brane}$, $\Gamma_{\rm bulk}$ & — & section 22.2 of 7+1.nb & radion decay rates into brane fields and into the bulk \\
$Q$, $Q'$ & — & $V=QQ'/(4\pi r)$ & electric charges on the brane \\*
\end{longtable}}
\begin{align*}
3H_a^2&=8\pi G_4\rho_4+\Lambda_4,\qquad H_a'=-\frac\kappa2(\varepsilon+p_a),&&(5.1) \\
8\pi G_4&=\frac\kappa{V_\perp}=\frac1{M_{\rm Pl}^2},\qquad \Lambda_4=\Lambda+\kappa U_0+\frac{\kappa f_W}{V_\perp}, \\
t_0&=\frac1{H_0}\int_0^1\frac{da}{a\sqrt{\Omega_ra^{-4}+\Omega_ma^{-3}+\Omega_\Lambda}},&&(5.2) \\
\Phi_8&=-\frac{5\kappa M}{32\pi^3r^5},\qquad \Phi_{\rm N}=-\frac{5\kappa M}{32\pi^3}\sum_{\vec n}\frac1{(r^2+\ell_\perp^2|\vec n|^2)^{5/2}},&&(5.3) \\
\Phi_{\rm N}&=-\frac53\frac{G_4M}r\sum_{\vec m}e^{-2\pi|\vec m|r/\ell_\perp}, \\
\beta&=\ln\frac b{b_0},\qquad U(b)=U_0+\frac{K_U}2\ln^2\frac b{b_0},\qquad m_r^2=\frac{\kappa K_U}{12},\qquad \phi=\sqrt{12}M_{\rm Pl}\beta,&&(5.4) \\
h&=e^{-4\beta}g_E,\qquad \mathcal L_{\phi W}=\frac{\phi\,T^\mu{}_\mu}{\sqrt3M_{\rm Pl}},\qquad V_\phi=-\frac23\frac{G_4MM'}re^{-m_rr}.
\end{align*}
\Needspace{12\baselineskip}\section*{6. The Higgs boson and the radion}\addcontentsline{toc}{section}{6. The Higgs boson and the radion}
{\footnotesize
\begin{longtable}{>{\raggedright\arraybackslash}p{\dimexpr 0.1700\linewidth-2\tabcolsep\relax}>{\raggedright\arraybackslash}p{\dimexpr 0.2200\linewidth-2\tabcolsep\relax}>{\raggedright\arraybackslash}p{\dimexpr 0.3100\linewidth-2\tabcolsep\relax}>{\raggedright\arraybackslash}p{\dimexpr 0.3000\linewidth-2\tabcolsep\relax}}
\toprule
\textbf{symbol} & \textbf{Mathematica} & \textbf{definition} & \textbf{meaning} \\
\midrule\endfirsthead
\toprule
\textbf{symbol} & \textbf{Mathematica} & \textbf{definition} & \textbf{meaning} \\
\midrule\endhead
\endfoot
\bottomrule\endlastfoot
$\Phi$, $\tilde\Phi$ & — & $\tilde\Phi=i\sigma_2\Phi^*$ & Higgs doublet; its conjugate \\
$\varphi_H$, $h$ & \texttt{rdPhE}, \texttt{rdHh} & $\Phi^\dagger\Phi=\varphi_H^2/2$, $\varphi_H=v+h$ & radial Higgs field; the Higgs field $h$ (not the metric $h_{\mu\nu}$) \\
$v$ & \texttt{vEW} & $(\sqrt2G_F)^{-1/2}$ & Higgs vacuum value \\
$\lambda$ & \texttt{lamH}, \texttt{rdLamH} & $m_H^2/(2v^2)$ & Higgs quartic coupling \\
$\xi$ & \texttt{rdXi} & (6.1) & curvature coupling of the Higgs ($\frac16$: conformal) \\
$R(h)$ & \texttt{rdRh} & scalar curvature of $h_{\mu\nu}$ & curvature of the brane metric \\
$\Omega$ & \texttt{rdOmg}, \texttt{rdOmE} & (6.2) & conformal factor of the brane metric \\
$\Box$ & — & $\partial_0^2-\nabla^2$ & flat wave operator \\
$\gamma$ & \texttt{rdGam} & $v/(\sqrt3M_{\rm Pl})$ & Higgs–radion mixing parameter \\
$m_{h0}$, $m_{r0}$ & \texttt{rdMh0}, \texttt{rdMr0} & $m_{h0}^2=2\lambda v^2$ & bare Higgs and radion masses \\
$h'$, $\phi'$, $Z$ & \texttt{rdSm}, \texttt{rdZ} & (6.3) & shifted fields; normalization \\
$\theta$, $h_m$, $\phi_m$ & \texttt{rdth}, \texttt{rdQ} & (6.3) & mixing angle; mass eigenstates \\*
\end{longtable}}
\begin{align*}
\mathcal L&=\sqrt{|h|}\Big[(D_\mu\Phi)^\dagger D^\mu\Phi-\lambda\Big(\Phi^\dagger\Phi-\frac{v^2}2\Big)^2+\xi R(h)\Phi^\dagger\Phi\Big]+\frac12(\partial\phi)^2-\frac12m_{r0}^2\phi^2,&&(6.1) \\
h_{\mu\nu}&=\Omega^2g_{E\mu\nu},\qquad \Omega=e^{-2\beta}=e^{-\phi/(\sqrt3M_{\rm Pl})},\qquad \sqrt{|h|}R(h)=-6\Omega\Box\Omega\ \ (g_E=\eta),&&(6.2) \\
h&=h'-\frac{6\xi\gamma}Z\phi',\qquad \phi=\frac{\phi'}Z,\qquad Z^2=1+6\xi\gamma^2(1-6\xi),&&(6.3) \\
h'&=\cos\theta\,h_m-\sin\theta\,\phi_m,\qquad \phi'=\sin\theta\,h_m+\cos\theta\,\phi_m,\qquad \tan2\theta=\frac{12\gamma\xi Zm_{h0}^2}{m_{r0}^2-m_{h0}^2(Z^2-36\xi^2\gamma^2)} .
\end{align*}
\Needspace{12\baselineskip}\section*{7. The Standard Model on the brane}\addcontentsline{toc}{section}{7. The Standard Model on the brane}
{\footnotesize
\begin{longtable}{>{\raggedright\arraybackslash}p{\dimexpr 0.1700\linewidth-2\tabcolsep\relax}>{\raggedright\arraybackslash}p{\dimexpr 0.2200\linewidth-2\tabcolsep\relax}>{\raggedright\arraybackslash}p{\dimexpr 0.3100\linewidth-2\tabcolsep\relax}>{\raggedright\arraybackslash}p{\dimexpr 0.3000\linewidth-2\tabcolsep\relax}}
\toprule
\textbf{symbol} & \textbf{Mathematica} & \textbf{definition} & \textbf{meaning} \\
\midrule\endfirsthead
\toprule
\textbf{symbol} & \textbf{Mathematica} & \textbf{definition} & \textbf{meaning} \\
\midrule\endhead
\endfoot
\bottomrule\endlastfoot
$G^a_\mu$, $W^I_\mu$, $B_\mu$ & — & gauge fields & $SU(3)$ ($a=1..8$), $SU(2)$ ($I=1,2,3$), $U(1)_Y$ \\
$g_s$, $g$, $g'$ & \texttt{gEW}, \texttt{gpEW} & $g=2m_W/v$, $g'=g\tan\theta_W$ & gauge couplings \\
$Q$, $T_3$, $Y$ & — & $Q=T_3+Y$ & electric charge, weak isospin, hypercharge \\
$Q_L$, $u_R$, $d_R$, $L_L$, $e_R$, $\nu_R$ & — & table of section 8 of the document & the Weyl fields of one generation \\
$g_*$, $g_{*s}$ & — & $28+\frac78\cdot90=106.75$ & relativistic degrees of freedom (energy, entropy) \\
$\gamma^\mu(x)$ & — & $\{\gamma^\mu,\gamma^\nu\}=2h^{\mu\nu}$ & Dirac matrices of the brane metric \\
$D_\mu$, $\nabla_\mu$ & — & (7.1) & gauge-covariant derivative; covariant derivative of $h$ \\
$T^a$, $\tau^I$, $f^{abc}$ & — & (7.1) & $SU(3)$ generators, $\tau^I=\sigma_I$, structure constants \\
$Y_u$, $Y_d$, $Y_e$, $Y_\nu$ & — & $3\times3$ matrices & Yukawa matrices \\
$\mu^2$ & — & $-\lambda v^2$ & Higgs mass parameter \\
$G_F$, $m_W$, $m_Z$, $m_H$ & \texttt{gFermi}, \texttt{mWb}, \texttt{mZb}, \texttt{mHb} & PDG 2025 & inputs \\
$\theta_W$, $e$ & \texttt{sW2}, \texttt{eEW} & $\sin^2\theta_W=1-m_W^2/m_Z^2$, $e=g\sin\theta_W$ & weak mixing angle; electric charge \\*
\end{longtable}}
\begin{align*}
D_\mu&=\nabla_\mu-ig_sG^a_\mu T^a-igW^I_\mu\frac{\tau^I}2-ig'YB_\mu,\qquad G^a_{\mu\nu}=\partial_\mu G^a_\nu-\partial_\nu G^a_\mu+g_sf^{abc}G^b_\mu G^c_\nu,&&(7.1) \\
\mathcal L_{\rm SM}&=-\frac14G^a_{\mu\nu}G^{a\mu\nu}-\frac14W^I_{\mu\nu}W^{I\mu\nu}-\frac14B_{\mu\nu}B^{\mu\nu}+(D_\mu\Phi)^\dagger D^\mu\Phi-\mu^2\Phi^\dagger\Phi-\lambda(\Phi^\dagger\Phi)^2 \\
&\quad+\sum_f\bar fi\gamma^\mu D_\mu f-[\text{Yukawa terms}+\text{h.c.}].
\end{align*}
\Needspace{12\baselineskip}\section*{8. dirac16complex at the brane}\addcontentsline{toc}{section}{8. dirac16complex at the brane}
{\footnotesize
\begin{longtable}{>{\raggedright\arraybackslash}p{\dimexpr 0.1700\linewidth-2\tabcolsep\relax}>{\raggedright\arraybackslash}p{\dimexpr 0.2200\linewidth-2\tabcolsep\relax}>{\raggedright\arraybackslash}p{\dimexpr 0.3100\linewidth-2\tabcolsep\relax}>{\raggedright\arraybackslash}p{\dimexpr 0.3000\linewidth-2\tabcolsep\relax}}
\toprule
\textbf{symbol} & \textbf{Mathematica} & \textbf{definition} & \textbf{meaning} \\
\midrule\endfirsthead
\toprule
\textbf{symbol} & \textbf{Mathematica} & \textbf{definition} & \textbf{meaning} \\
\midrule\endhead
\endfoot
\bottomrule\endlastfoot
$\Gamma_A$, $\Gamma^A$ & \texttt{Γ17[A]}, \texttt{Γ17up[A]} & (8.1) & real $16\times16$ Dirac matrices of 7+1 \\
$\sigma_i$ & \texttt{PauliMatrix[i]} & (8.1) & Pauli matrices \\
$\gamma_\mu$, $\gamma_5$ & \texttt{g4D}, \texttt{g5D} & (8.1) & $4\times4$ Dirac matrices of the brane \\
$\rho_k$ & \texttt{rhoK} & (8.1) & anti-Hermitian $4\times4$ matrices of the transverse directions \\
$R_A$ & \texttt{rBr} & (8.1) & brane-adapted $16\times16$ representation \\
$U$ & \texttt{uInt}, \texttt{tbU8} ($2\sqrt2U$) & (8.2) & unitary intertwiner \\
$\Gamma^{(K)}$, $R^{(K)}$, $s$ & — & (8.2) & Clifford products; a matrix unit \\
$\zeta$, $\psi_a$, $e_a$ & — & (8.2) & brane-adapted field; the four brane spinors; unit vectors of $\mathbb C^4$ \\
$u$ & \texttt{tbu} & $(1,2,0,-1)$ & the brane spinor of Example 9.2 \\
$\rho_{4567}$, $P_s$ & \texttt{rho4567}, \texttt{tbPs} & (8.3) & transverse chirality; its projectors \\
$\chi$, $P_L$ & \texttt{PL17} & (8.3) & eight-dimensional chirality; left projector \\
$X$ (operator), $m_{\vec n}$, $\chi_{\vec n}$ & \texttt{xKK} & (8.4) & Kaluza–Klein operator; masses; four-dimensional fields \\
$r_4(k)$, $\sigma(k)$ & \texttt{rIIr4}, \texttt{rIIsig} & $8\sigma(k)-32\sigma(k/4)$ & lattice multiplicity; divisor sum \\*
\end{longtable}}
\begin{align*}
\Gamma_A\Gamma_B+\Gamma_B\Gamma_A&=2\eta_{AB}1_{16},\qquad \gamma_0=\mathrm{diag}(1,1,-1,-1),&&(8.1) \\
\gamma_i&=\begin{pmatrix}0&\sigma_i\\-\sigma_i&0\end{pmatrix},\qquad \gamma_5=i\gamma_0\gamma_1\gamma_2\gamma_3, \\
\rho_4&=i\gamma_0,\quad \rho_{5,6,7}=-\gamma_{1,2,3},\qquad R_\mu=\gamma_\mu\otimes1_4,\quad R_k=\gamma_5\otimes\rho_k, \\
\Gamma_AU&=UR_A,\qquad U\propto\sum_K\Gamma^{(K)}s\,(R^{(K)})^{-1},\qquad \psi=U\zeta,\qquad \zeta=\sum_{a=1}^4\psi_a\otimes e_a,&&(8.2) \\
\rho_{4567}&=\rho_4\rho_5\rho_6\rho_7=-\gamma_5,\qquad P_s=\tfrac12(1_4+s\rho_{4567}),&&(8.3) \\
\chi&=i\Gamma_0\cdots\Gamma_7=U(\gamma_5\otimes\rho_{4567})U^\dagger,\qquad P_L=\tfrac12(1-\chi), \\
X&=\gamma_5\otimes\frac{\vec\rho\cdot\vec q}{b_0},\qquad X^2=-\frac{|\vec q|^2}{b_0^2},\qquad m_{\vec n}=\sqrt{m^2+\frac{4\pi^2|\vec n|^2}{\ell_\perp^2}},&&(8.4) \\
\psi(x,0)&=\frac1{\sqrt{V_\perp}}\sum_{\vec n}\chi_{\vec n}(x).
\end{align*}
\Needspace{12\baselineskip}\section*{9. The portals}\addcontentsline{toc}{section}{9. The portals}
{\footnotesize
\begin{longtable}{>{\raggedright\arraybackslash}p{\dimexpr 0.1700\linewidth-2\tabcolsep\relax}>{\raggedright\arraybackslash}p{\dimexpr 0.2200\linewidth-2\tabcolsep\relax}>{\raggedright\arraybackslash}p{\dimexpr 0.3100\linewidth-2\tabcolsep\relax}>{\raggedright\arraybackslash}p{\dimexpr 0.3000\linewidth-2\tabcolsep\relax}}
\toprule
\textbf{symbol} & \textbf{Mathematica} & \textbf{definition} & \textbf{meaning} \\
\midrule\endfirsthead
\toprule
\textbf{symbol} & \textbf{Mathematica} & \textbf{definition} & \textbf{meaning} \\
\midrule\endhead
\endfoot
\bottomrule\endlastfoot
$y^L_{is}$, $y^R_{is}$, $y$ & \texttt{cyS} ($c_y$) & (9.1) & Yukawa (neutrino) portal couplings \\
$\psi^{(L)}_s$, $\psi^{(R)}_s$, $P_R$ & — & $P_R=\frac12(1+\gamma_5)$ & brane-right-handed parts of the chiral halves \\
$\mu_{is}$, $\mu_P$ & \texttt{dkMuP} & (9.1) & portal Dirac masses \\
$G^f_P$, $G_P$ & \texttt{cVS} ($c_V$) & (9.1) & vector portal \\
$\lambda_P$ & \texttt{cPS} ($c_P$) & (9.1) & Higgs portal \\
$\phi_{\vec n}$ & \texttt{tbPhi[n]} & $e^{i\vec q_{\vec n}\cdot\vec y}/\sqrt{V_\perp}$ & normalized modes of the torus \\
$c_P$, $c_y$, $c_V$ & \texttt{cPS}, \texttt{cyS}, \texttt{cVS} & (9.2) & couplings in units of $M_8$ \\
$\lambda_{\rm eff}$, $y_{\rm eff}$, $G_{\rm eff}$ & \texttt{dkLam}, \texttt{dkYeff} & (9.2) & four-dimensional couplings of one level \\
$\mathcal M$, $m_D$, $N$ & — & (9.3) & neutrino mass matrix; Dirac masses; light mixing block \\
$U_{\rm PMNS}$, $m_1,m_2,m_3$ & — & section 17 of 7+1.nb & lepton mixing matrix; light neutrino masses \\*
\end{longtable}}
\begin{align*}
\mathcal L_{\rm portal}&\supset-\sum_{i,s}\big[y^L_{is}\bar L_i\tilde\Phi P_R\psi^{(L)}_s+y^R_{is}\bar L_i\tilde\Phi P_R\psi^{(R)}_s\big]+\text{h.c.}&&(9.1) \\
&\quad-G^f_P(\bar\psi\Gamma^\mu\psi)(\bar f\gamma_\mu f)-\lambda_P\Phi^\dagger\Phi\,\bar\psi\psi,\qquad \mu_{is}=\frac{y_{is}v}{\sqrt2\sqrt{V_\perp}}, \\
\lambda_P&=\frac{c_P}{M_8^5},\qquad y=\frac{c_y}{M_8^2},\qquad G_P=\frac{c_V}{M_8^6},&&(9.2) \\
\lambda_{\rm eff}&=\frac{c_PM_8}{M_{\rm Pl}^2},\qquad y_{\rm eff}=\frac{c_yM_8}{M_{\rm Pl}},\qquad G_{\rm eff}=\frac{c_V}{M_{\rm Pl}^2},\qquad \mu_P=\frac{y_{\rm eff}v}{\sqrt2}, \\
\mathcal M&=\begin{pmatrix}m_D&\mu\\0&m1_4\end{pmatrix},\qquad m_D=U_{\rm PMNS}\mathrm{diag}(m_1,m_2,m_3),\qquad 1-NN^\dagger\simeq\frac{\mu\mu^\dagger}{m^2}.&&(9.3)
\end{align*}
\Needspace{12\baselineskip}\section*{10. The three regimes}\addcontentsline{toc}{section}{10. The three regimes}
{\footnotesize
\begin{longtable}{>{\raggedright\arraybackslash}p{\dimexpr 0.1700\linewidth-2\tabcolsep\relax}>{\raggedright\arraybackslash}p{\dimexpr 0.2200\linewidth-2\tabcolsep\relax}>{\raggedright\arraybackslash}p{\dimexpr 0.3100\linewidth-2\tabcolsep\relax}>{\raggedright\arraybackslash}p{\dimexpr 0.3000\linewidth-2\tabcolsep\relax}}
\toprule
\textbf{symbol} & \textbf{Mathematica} & \textbf{definition} & \textbf{meaning} \\
\midrule\endfirsthead
\toprule
\textbf{symbol} & \textbf{Mathematica} & \textbf{definition} & \textbf{meaning} \\
\midrule\endhead
\endfoot
\bottomrule\endlastfoot
$T$ & — & temperature & temperature of the brane plasma \\
$T_*$ & \texttt{tbTstar} & $(M_8^6/M_{\rm Pl})^{1/5}$ & normalcy temperature (ESTIMATE) \\
$s$ (entropy) & — & $\frac{2\pi^2}{45}g_{*s}T^3$ & entropy density \\
$t$, $z$ & — & cosmic time; redshift & — \\
$H_\infty$ & \texttt{rIIIH} & $H_0\Omega_\Lambda^{1/2}$ & final de Sitter rate \\*
\end{longtable}}
\Needspace{12\baselineskip}\section*{11. The brane on a 4-sphere}\addcontentsline{toc}{section}{11. The brane on a 4-sphere}
{\footnotesize
\begin{longtable}{>{\raggedright\arraybackslash}p{\dimexpr 0.1700\linewidth-2\tabcolsep\relax}>{\raggedright\arraybackslash}p{\dimexpr 0.2200\linewidth-2\tabcolsep\relax}>{\raggedright\arraybackslash}p{\dimexpr 0.3100\linewidth-2\tabcolsep\relax}>{\raggedright\arraybackslash}p{\dimexpr 0.3000\linewidth-2\tabcolsep\relax}}
\toprule
\textbf{symbol} & \textbf{Mathematica} & \textbf{definition} & \textbf{meaning} \\
\midrule\endfirsthead
\toprule
\textbf{symbol} & \textbf{Mathematica} & \textbf{definition} & \textbf{meaning} \\
\midrule\endhead
\endfoot
\bottomrule\endlastfoot
$S^4$, $r$, $r_0$ & \texttt{rdR0} & (11.1) & the 4-sphere; its radius; the reference radius \\
$\hat\gamma$ & \texttt{rdGam4} & (11.1) & metric of the unit 4-sphere \\
$X_1,\dots,X_5$ & \texttt{rdEmb} & (11.1) & embedding in $\mathbb R^5$; the north pole $(0,0,0,0,r)$ \\
$L_{ab}$ & — & $X_a\partial_b-X_b\partial_a$ & the 10 rotations ($SO(5)$) \\
$l$, $\mu_l$, $d_l$ & \texttt{rdDl} & (11.2) & degree, mass and multiplicity of a level \\
$Z_l$, $C_l^{(3/2)}$ & \texttt{rdZon} & (11.2) & zonal harmonic; Gegenbauer polynomial \\
$r_{S^4}$ & \texttt{rdRS4} & $(3V_\perp/(8\pi^2))^{1/4}$ & radius of a sphere of volume $V_\perp$ \\
$\tilde\gamma_a$, $e^a$, $\varepsilon_s$ & \texttt{rdGt}, \texttt{rdHf}, \texttt{rdKS} & (11.3) & Dirac matrices, frame and Killing spinors of $S^4$ \\
$\lambda$ (Dirac), $R_{\rm int}$ & — & (11.3) & Dirac eigenvalues; $R_{\rm int}=12/r^2$ \\
$U_1$ & \texttt{rdU1} & $U=U_0+U_1\beta+\frac{K_U}2\beta^2$ & slope of the stabilizing energy (sphere) \\*
\end{longtable}}
\begin{align*}
g_{kl}&=-r^2\hat\gamma_{kl},\qquad \hat\gamma=\mathrm{diag}(1,\sin^2x_4,\sin^2x_4\sin^2x_5,\sin^2x_4\sin^2x_5\sin^2x_6),\qquad V=\frac{8\pi^2}3r^4,&&(11.1) \\
X_5&=r\cos x_4,\qquad X_4=r\sin x_4\cos x_5,\qquad X_3=r\sin x_4\sin x_5\cos x_6, \\
(X_1,X_2)&=r\sin x_4\sin x_5\sin x_6\,(\cos x_7,\sin x_7), \\
\mu_l&=\frac{\sqrt{l(l+3)}}r,\qquad d_l=\frac{(2l+3)(l+1)(l+2)}6,&&(11.2) \\
Z_l&=\sqrt{\frac{d_l}V}\frac{C_l^{(3/2)}(\cos x_4)}{C_l^{(3/2)}(1)},\qquad C_l^{(3/2)}(1)=\frac{(l+1)(l+2)}2, \\
\tilde\gamma_a&=i\rho_{a+3},\qquad e^a=(r,r\sin x_4,\dots)dx^{a+3},\qquad \nabla_i\varepsilon_s=s\frac i{2r}e_i{}^a\tilde\gamma_a\varepsilon_s,&&(11.3) \\
\lambda&=\pm\frac{i(l+2)}r,\qquad |\lambda|^2\ge\frac{nR_{\rm int}}{4(n-1)}=\frac4{r^2}.
\end{align*}
\Needspace{12\baselineskip}\section*{12. Warped branes (Randall–Sundrum)}\addcontentsline{toc}{section}{12. Warped branes (Randall–Sundrum)}
{\footnotesize
\begin{longtable}{>{\raggedright\arraybackslash}p{\dimexpr 0.1700\linewidth-2\tabcolsep\relax}>{\raggedright\arraybackslash}p{\dimexpr 0.2200\linewidth-2\tabcolsep\relax}>{\raggedright\arraybackslash}p{\dimexpr 0.3100\linewidth-2\tabcolsep\relax}>{\raggedright\arraybackslash}p{\dimexpr 0.3000\linewidth-2\tabcolsep\relax}}
\toprule
\textbf{symbol} & \textbf{Mathematica} & \textbf{definition} & \textbf{meaning} \\
\midrule\endfirsthead
\toprule
\textbf{symbol} & \textbf{Mathematica} & \textbf{definition} & \textbf{meaning} \\
\midrule\endhead
\endfoot
\bottomrule\endlastfoot
$A(y)$, $k$, $y$, $R$ & \texttt{rdA}, \texttt{rdk}, \texttt{rdy}, \texttt{rdRR} & (12.1) & warp function, AdS scale, fifth coordinate, radius \\
$f_{\rm UV}$, $f_{\rm IR}$ & \texttt{rdfUV}, \texttt{rdfIR} & (12.1) & brane tensions \\
$K^\mu{}_\nu$, $S_{\mu\nu}$ & \texttt{rdKmix}, \texttt{rdKjump} & (12.1) & extrinsic curvature; brane energy–momentum \\
$M$ (RS) & — & $\kappa=1/(4M^3)$ & five-dimensional scale of Randall and Sundrum \\
$T(x)$, $\mu$, $f$, $\langle\mu\rangle$, $\delta\mu$ & \texttt{rdT}, \texttt{rdmu}, \texttt{rdfm}, \texttt{rdmu0} & (12.2) & radius field; canonical radion; its vacuum value and fluctuation \\
$v_A$, $v_B$, $m$, $\epsilon$ & \texttt{rdvA}, \texttt{rdvB}, \texttt{rdMg}, \texttt{rdeps} & (12.2) & Goldberger–Wise brane values, mass, small parameter \\
$\vec z$ & \texttt{rdz1} … \texttt{rdz3} & — & the three unwarped extra directions of the eight-dimensional test \\*
\end{longtable}}
\begin{align*}
&ds^2=e^{-2A(y)}\eta-dy^2,\qquad G^\mu{}_\nu=3A''-6A'^2,\qquad G^y{}_y=-6A'^2,&&(12.1) \\
&f_{\rm UV}=-f_{\rm IR}=\frac{6k}\kappa,\qquad K^\mu{}_\nu=\frac12g^{\mu\alpha}\partial_yg_{\alpha\nu}=-A'\delta^\mu_\nu, \\
&\mu=fe^{-k\pi T},\qquad f^2=\frac6{\kappa k},\qquad M_{\rm Pl}^2=\frac{1-e^{-2\pi kR}}{\kappa k},\qquad \mu=\langle\mu\rangle+\delta\mu,&&(12.2) \\
&m_r^2=\frac43\kappa k^2\epsilon^2v_B^2\Big(\frac{\langle\mu\rangle}f\Big)^2,\qquad \epsilon=\frac{m^2}{4k^2}.
\end{align*}
\Needspace{12\baselineskip}\section*{13. Constants and data}\addcontentsline{toc}{section}{13. Constants and data}
{\footnotesize
\begin{longtable}{>{\raggedright\arraybackslash}p{\dimexpr 0.1700\linewidth-2\tabcolsep\relax}>{\raggedright\arraybackslash}p{\dimexpr 0.1700\linewidth-2\tabcolsep\relax}>{\raggedright\arraybackslash}p{\dimexpr 0.3000\linewidth-2\tabcolsep\relax}>{\raggedright\arraybackslash}p{\dimexpr 0.3600\linewidth-2\tabcolsep\relax}}
\toprule
\textbf{symbol} & \textbf{Mathematica} & \textbf{value} & \textbf{meaning and source} \\
\midrule\endfirsthead
\toprule
\textbf{symbol} & \textbf{Mathematica} & \textbf{value} & \textbf{meaning and source} \\
\midrule\endhead
\endfoot
\bottomrule\endlastfoot
$M_8$ & \texttt{rIIM8} & $10^4$ GeV & choice of sections 22–24 \\
$M_{\rm Pl}$ & \texttt{rIIMPl} & $2.435\times10^{18}$ GeV & reduced Planck mass \\
$m_r$ & — & 10 GeV & radion mass (choice of sections 22–24) \\
$m_h$, $m_H$ & \texttt{mHb} & 125.20 GeV & Higgs mass (PDG 2025) \\
$\hbar c$ & \texttt{tbHbarCm} & $1.973269804\times10^{-16}$ GeV m & CODATA 2018 \\
$H_0$, $\Omega_m$, $\Omega_r$ & — & 67.4 km/s/Mpc, 0.315, $9.14\times10^{-5}$ & Planck 2018 \\
$\langle\mu\rangle$ & — & 3 TeV & the warped example of the author’s note \\*
\end{longtable}}
\Needspace{12\baselineskip}\section*{14. Symbols that look alike}\addcontentsline{toc}{section}{14. Symbols that look alike}
\begin{itemize}
\item $h$: the induced metric $h_{\mu\nu}$ (always with indices or as a matrix), the Higgs field $h$ of section 6 (and $h'$, $h_m$).
\item $\Phi$: the Higgs doublet $\Phi$; $\Phi_8$, $\Phi_{\rm N}$ are gravitational potentials; $\varphi_H$ the radial Higgs field.
\item $\phi$: the canonical radion $\phi$ (and $\phi'$, $\phi_m$) and the Fourier modes $\phi_{\vec n}$ (always with an index).
\item $\chi$: the eight-dimensional chirality $\chi$ and the four-dimensional fields $\chi_{\vec n}$ (with an index).
\item $\lambda$: the Higgs quartic $\lambda$, the Higgs portal $\lambda_P$, $\lambda_{\rm eff}$, the Dirac eigenvalue $\lambda$ on $S^4$, and the index of the Christoffel symbol $\{\lambda\,\mu\nu\}$.
\item $\mu$: the index $\mu=0,\dots,3$; the portal masses $\mu_{is}$, $\mu_P$, $\mu$; the Kaluza–Klein masses $\mu_l$; the warped radion $\mu$; $\mu^2$ of the Higgs potential.
\item $\gamma$: the brane Dirac matrices $\gamma_\mu$, $\gamma^\mu(x)$, $\gamma_5$; the mixing parameter $\gamma$; $\hat\gamma$ and $\tilde\gamma_a$ of the sphere.
\item $\rho$: the densities $\rho_4$, $\rho_\Lambda$; the matrices $\rho_k$, $\rho_{4567}$.
\item $\sigma$: the Pauli matrices $\sigma_i$ and the divisor sum $\sigma(k)$.
\item $r$: the distance $r$ between two masses; the radius $r$ of the sphere ($r_0$, $r_{S^4}$); $r_4(k)$.
\item $R$: the scalar curvature $R$, $R(h)$, $R_{\rm int}$; the Randall–Sundrum radius $R$; the matrices $R_A$, $R^{(K)}$.
\item $K$: the extrinsic curvatures $K^{(k)}_{\mu\nu}$, $K^\mu{}_\nu$; the stiffness $K_U$; the index $K$ of a Clifford product.
\item $S$: the action and its parts; the brane energy–momentum $S_{\mu\nu}$ of section 12; the sphere $S^4$; $s$ is a sign or the proper time or the entropy.
\item $M$: the bulk $M$; masses $M$, $M'$; the scales $M_8$, $M_{\rm Pl}$, $M$ (RS); the mass matrix $\mathcal M$.
\item $U$: the intertwiner $U$; the stabilizing potential $U(b)$ with $U_0$, $U_1$; $U_{\rm PMNS}$.
\item $W$: the brane $W$; the comoving energy $W(a,b)$; the gauge fields $W^I_\mu$.
\item $Q$: electric charges $Q$, $Q'$ and the quark doublet $Q_L$.
\item $f$: the brane tension $f_W$ and $f_{\rm UV}$, $f_{\rm IR}$; the normalization $f$ of the warped radion; Standard Model fermions $f$; $f^{abc}$.
\item $\epsilon$, $\varepsilon$: the Levi-Civita symbol $\epsilon_{ijk}$ and the Goldberger–Wise parameter $\epsilon$; the energy density $\varepsilon$ and the Killing spinors $\varepsilon_s$.
\end{itemize}
\Needspace{12\baselineskip}\section*{15. Mathematica names}\addcontentsline{toc}{section}{15. Mathematica names}
\Needspace{8\baselineskip}\subsection*{15.1 Names of section 25 used by the document}\addcontentsline{toc}{subsection}{15.1 Names of section 25 used by the document}
{\footnotesize
\begin{longtable}{>{\raggedright\arraybackslash}p{\dimexpr 0.4200\linewidth-2\tabcolsep\relax}>{\raggedright\arraybackslash}p{\dimexpr 0.5800\linewidth-2\tabcolsep\relax}}
\toprule
\textbf{name} & \textbf{meaning} \\
\midrule\endfirsthead
\toprule
\textbf{name} & \textbf{meaning} \\
\midrule\endhead
\endfoot
\bottomrule\endlastfoot
\texttt{tbY}, \texttt{tbLc}, \texttt{tbB}, \texttt{tbV}, \texttt{tbLp}, \texttt{tbHbarCm} & $\vec y$; $L_c$; $b$; $V_\perp$; $\ell_\perp$; $\hbar c$ \\
\texttt{tbJ}, \texttt{tbOnW}, \texttt{tbInduced}, \texttt{tbPt}, \texttt{tbChrW} & the Jacobian $J$; the rule $\vec y\to0$; $h=J^{\mathsf T}gJ$; the point $x_0=\frac12$, $x_4=0$; Christoffel symbols at the brane \\
\texttt{tbQ}, \texttt{tbPhi}, \texttt{tbVtx}, \texttt{tbRho} & $\vec q_{\vec n}$; $\phi_{\vec n}$; overlap of three modes; $\rho_4$ \\
\texttt{tbKap}, \texttt{tbGN}, \texttt{tbr}, \texttt{tbMP}, \texttt{tbS}, \texttt{tbSkk} & symbols $\kappa$, $G$, $r$, $M_{\rm Pl}$; the image and Kaluza–Klein sums \\
\texttt{tbSPermQ}, \texttt{tbU8}, \texttt{tbUd}, \texttt{tbKp}, \texttt{tbZ16} & signed-permutation test; $2\sqrt2U$; $U^\dagger$; \texttt{KroneckerProduct}; the $16\times16$ zero matrix \\
\texttt{tbPs}, \texttt{tbEven} & $P_s$; the even permutations of $(4,5,6,7)$ \\
\texttt{tbu}, \texttt{tbPsi}, \texttt{tbBil} & $u=(1,2,0,-1)$; $\psi=U(u\otimes e_1)$; $\bar\psi\Gamma^\mu\psi$ \\
\texttt{tbfW}, \texttt{tbu0} … \texttt{tbu3}, \texttt{tbAa}, \texttt{tbBb}, \texttt{tbr0}, \texttt{tbth} & $f_W$; components of $h^{\mu\nu}$; arguments $a$, $b$ of $W(a,b)$; a radiation constant; $\theta$ \\
\texttt{tbT}, \texttt{tbM8}, \texttt{tbTstar}, \texttt{tby} & symbols $T$, $M_8$; $T_*$; $|\vec y|$ \\*
\end{longtable}}
\Needspace{8\baselineskip}\subsection*{15.2 Names of section 26 used by the document}\addcontentsline{toc}{subsection}{15.2 Names of section 26 used by the document}
{\footnotesize
\begin{longtable}{>{\raggedright\arraybackslash}p{\dimexpr 0.4200\linewidth-2\tabcolsep\relax}>{\raggedright\arraybackslash}p{\dimexpr 0.5800\linewidth-2\tabcolsep\relax}}
\toprule
\textbf{name} & \textbf{meaning} \\
\midrule\endfirsthead
\toprule
\textbf{name} & \textbf{meaning} \\
\midrule\endhead
\endfoot
\bottomrule\endlastfoot
\texttt{rdBt}, \texttt{rdBe}, \texttt{rdR0}, \texttt{rdGam4}, \texttt{rdEmb}, \texttt{rdDl}, \texttt{rdZon}, \texttt{rdRS4} & $\beta(x_0)$, $\beta$, $r_0$, $\hat\gamma$, $X_a$, $d_l$, $C_l^{(3/2)}(\cos x_4)$, $r_{S^4}$ \\
\texttt{rdGt}, \texttt{rdHf}, \texttt{rdKS}, \texttt{rdU1} & $\tilde\gamma_a$; frame lengths; Killing spinors; $U_1$ \\
\texttt{rdA}, \texttt{rdk}, \texttt{rdy}, \texttt{rdRR}, \texttt{rdfUV}, \texttt{rdfIR}, \texttt{rdKmix}, \texttt{rdKjump} & RS1: $A(y)$, $k$, $y$, $R$, tensions, $K^\mu{}_\nu$ and its jump \\
\texttt{rdT}, \texttt{rdmu}, \texttt{rdfm}, \texttt{rdmu0}, \texttt{rdvA}, \texttt{rdvB}, \texttt{rdMg}, \texttt{rdeps}, \texttt{rdz1} … \texttt{rdz3} & the warped radion and its Goldberger–Wise stabilization; the unwarped directions \\
\texttt{rdPhE}, \texttt{rdHh}, \texttt{rdLamH}, \texttt{rdXi}, \texttt{rdRh}, \texttt{rdOmg}, \texttt{rdOmE}, \texttt{rdGam} & $\varphi_H$, $h$, $\lambda$, $\xi$, $R(h)$, $\Omega$, $\gamma$ \\
\texttt{rdMh0}, \texttt{rdMr0}, \texttt{rdSm}, \texttt{rdZ}, \texttt{rdth}, \texttt{rdQ} & $m_{h0}$, $m_{r0}$, the shift, $Z$, $\theta$, the rotation \\*
\end{longtable}}
\Needspace{8\baselineskip}\subsection*{15.3 Names of the earlier sections}\addcontentsline{toc}{subsection}{15.3 Names of the earlier sections}
{\footnotesize
\begin{longtable}{>{\raggedright\arraybackslash}p{\dimexpr 0.4200\linewidth-2\tabcolsep\relax}>{\raggedright\arraybackslash}p{\dimexpr 0.1000\linewidth-2\tabcolsep\relax}>{\raggedright\arraybackslash}p{\dimexpr 0.4800\linewidth-2\tabcolsep\relax}}
\toprule
\textbf{name} & \textbf{section} & \textbf{meaning} \\
\midrule\endfirsthead
\toprule
\textbf{name} & \textbf{section} & \textbf{meaning} \\
\midrule\endhead
\endfoot
\bottomrule\endlastfoot
\texttt{X}, \texttt{ID4}, \texttt{ID16} & 1 & coordinates; unit matrices \\
\texttt{Γ17[A]}, \texttt{Γ17up[A]}, \texttt{η17} & 3–5 & Dirac matrices of 7+1, raised; $\eta$ \\
\texttt{gCos}, \texttt{a3[x0]}, \texttt{b4[x0]}, \texttt{cvCos}, \texttt{gPoly}, \texttt{nP}, \texttt{aP}, \texttt{bP}, \texttt{toH}, \texttt{Ha}, \texttt{Hb} & 11 & the cosmological and inhomogeneous metrics, curvature, expansion rates \\
\texttt{g4D}, \texttt{g5D}, \texttt{rhoK}, \texttt{rBr}, \texttt{uInt}, \texttt{rho4567}, \texttt{PL17} & 17.2 & $\gamma_\mu$, $\gamma_5$, $\rho_k$, $R_A$, $U$, $\rho_{4567}$, $P_L$ \\
\texttt{vEW}, \texttt{gFermi}, \texttt{mWb}, \texttt{mZb}, \texttt{mHb}, \texttt{gEW}, \texttt{gpEW}, \texttt{sW2}, \texttt{eEW}, \texttt{lamH} & 17.3 & electroweak constants \\
\texttt{xKK} & 18.1 & the Kaluza–Klein operator $X$ \\
\texttt{feStress}, \texttt{uPot}, \texttt{u0B}, \texttt{kU}, \texttt{lamb}, \texttt{fWb}, \texttt{alf} & 20 & stress of a comoving energy; stabilizing potential; $\Lambda$; $f_W$; Lovelock couplings \\
\texttt{rIp1}, \texttt{rIq1} & 21.1 & Kasner exponents \\
\texttt{rIIr4}, \texttt{rIIsig}, \texttt{rIIVp}, \texttt{rIIM8}, \texttt{rIIMPl}, \texttt{rIIH0g}, \texttt{rIIOL} & 22 & lattice sums; $V_\perp$, $M_8$, $M_{\rm Pl}$; $H_0$, $\Omega_\Lambda$ \\
\texttt{rIIIH} & 23 & $H_\infty$ \\
\texttt{cPS}, \texttt{cyS}, \texttt{cVS}, \texttt{dkLam}, \texttt{dkYeff}, \texttt{dkMuP}, \texttt{dkRhoL} & 24 & portal couplings; $\lambda_{\rm eff}$, $y_{\rm eff}$, $\mu_P$; $\rho_\Lambda$ \\*
\end{longtable}}
\end{document}
